% Generated draft: requires mathematical and source review.
\documentclass[11pt]{article}
\usepackage[margin=25mm]{geometry}
\usepackage{fontspec}
\setmainfont{Latin Modern Roman}
\usepackage{amsmath,amssymb,mathtools,unicode-math}
\setmathfont{Latin Modern Math}
\usepackage{enumitem,needspace,xurl,hyperref}
\setlength{\parindent}{0pt}
\setlength{\parskip}{5pt}
\setlength{\emergencystretch}{3em}
\newcommand{\N}{\mathbb{N}}
\newcommand{\Z}{\mathbb{Z}}
\newcommand{\Q}{\mathbb{Q}}
\newcommand{\R}{\mathbb{R}}
\newcommand{\C}{\mathbb{C}}
\newcommand{\F}{\mathbb{F}}
\newcommand{\A}{\mathbb{A}}
\newcommand{\PP}{\mathbb{P}}
\newcommand{\E}{\mathbb{E}}
\newcommand{\eps}{\varepsilon}
\newcommand{\dd}{\,\mathrm{d}}
\begin{document}
\section*{UnsolvedMath ID: 800669}
\textbf{Source problem code:} catalogue entry
\subsection{Definitions and mathematical statement}
Under the customary interpretation, for a topological space \(X\), let \(C_p(X)\) denote the set of continuous functions \(X\to\R\), equipped with the topology of pointwise convergence, and let \(\mathbf{0}\) be the constant zero function. The supplied record does not explicitly specify this convention or the hypotheses imposed on \(X\).

The record defines the \emph{monotonic sequence selection property} by modifying the sequence selection property, abbreviated SSP: for each \(n\in\N\), its input sequence \((f_k^n)_{k\in\N}\) must converge pointwise monotonically to \(\mathbf{0}\). In particular, the convergence requirement is
\[
\forall n\in\N\;\forall x\in X,\qquad \lim_{k\to\infty}f_k^n(x)=0.
\]
The monotonicity requirement is additional to this displayed limit. The supplied record does not specify its direction, whether that direction may depend on \(n\) or \(x\), or whether monotonicity refers to the function values or their absolute values.

The question is whether, for every space \(X\) in the intended class,
\[
C_p(X)\text{ has the monotonic sequence selection property}
\quad\Longrightarrow\quad
C_p(X)\text{ has the weak sequence selection property}.
\]
Neither SSP nor the weak sequence selection property is defined in the supplied record. Consequently, their selection requirements and quantifiers, and hence a complete self-contained formulation of this implication, remain unresolved. This is a statement-only draft requiring human mathematical review.
\subsection{Short English statement}
Does the version of SSP requiring each input sequence to converge pointwise monotonically to zero imply the weak sequence selection property for \(C_p(X)\)? The supplied excerpt omits the definitions of both underlying selection properties and the intended hypotheses on \(X\).
\subsection{Sources}
\begin{itemize}
\item Supplied local UnsolvedMath registry record: permanent ID \(800669\), exact source problem code \texttt{catalogue entry}; title: \emph{Selection Principles Milestone Problems --- Problem 42 ([CpHure, Problem 2])}. The registry attributes the problem to Tsaban, with proposed year 2003, and labels it \texttt{open}; that label is supplied metadata, not an independently verified status assessment.
\item Tsaban, \emph{Selection Principles in Mathematics: A Milestone of Open Problems} (2003), Problem 42. Supplied source URL: \url{https://arxiv.org/abs/math/0312182}. The record reports access on 2026-10-16 and extraction by \texttt{arxiv\_latex}; these are supplied provenance statements, not a claim of source access for this draft.
\item Supplied attribution/context: [CpHure, Problem 2]. The record provides no bibliographic expansion of the label \texttt{CpHure}.
\end{itemize}
\end{document}
